\documentclass[12pt]{article}
\usepackage[margin=1in]{geometry}
\usepackage[utf8]{inputenc}
\usepackage{latexsym, amsfonts, enumitem, graphicx, environ,enumitem, fancyhdr, color, amssymb, amsthm, amsmath, mathtools, tikz, nicematrix}
\usepackage{physics}
\usetikzlibrary{patterns}
\usetikzlibrary{decorations.pathreplacing}
\pagestyle{fancy}
\setlength{\headheight}{65pt}
\newenvironment{problem}[2][Problem]{\begin{trivlist}
    \item[\hskip \labelsep {\bfseries #1}\hskip \labelsep {\bfseries #2.}]}{\end{trivlist}}
\DeclareMathOperator{\Ima}{Im}

\usepackage{hyperref}
\usepackage[capitalize]{cleveref}

\newtheorem{thm}{Theorem}[section]
\newtheorem{lem}[thm]{Lemma}
\newtheorem{cl}[thm]{Claim}
\newtheorem{prop}[thm]{Proposition}
\newtheorem{cor}[thm]{Corollary}
\newtheorem{conj}[thm]{Conjecture}
\newtheorem{obstruction}[thm]{Obstruction}

\usepackage{mdframed}

\theoremstyle{definition}
\newtheorem{definition}[thm]{Definition}
\newtheorem{remark}[thm]{Remark}
\newtheorem{question}{Question}
\newtheorem{example}[thm]{Example}
\newtheorem{exercise}[thm]{Exercise}
\newtheorem{properties}[thm]{Properties}

\usepackage{tcolorbox}
\tcbuselibrary{theorems,skins,breakable}

\tcolorboxenvironment{example}{
enhanced jigsaw,colframe=cyan,interior hidden,
breakable,%before skip=10pt,after skip=10pt
}

\tcolorboxenvironment{thm}{
enhanced jigsaw,colframe=red,interior hidden,
breakable,%before skip=10pt,after skip=10pt
}

\tcolorboxenvironment{prop}{
enhanced jigsaw,colframe=red,interior hidden,
breakable,%before skip=10pt,after skip=10pt
}

\tcolorboxenvironment{properties}{
enhanced jigsaw,colframe=red,interior hidden,
breakable,%before skip=10pt,after skip=10pt
}

\tcolorboxenvironment{lem}{
enhanced jigsaw,colframe=red,interior hidden,
breakable,%before skip=10pt,after skip=10pt
}
\tcolorboxenvironment{cor}{
enhanced jigsaw,colframe=red,interior hidden,
breakable,%before skip=10pt,after skip=10pt
}
\tcolorboxenvironment{obstruction}{
enhanced jigsaw,colframe=red,interior hidden,
breakable,%before skip=10pt,after skip=10pt
}

\tcolorboxenvironment{proof}{
enhanced jigsaw, frame hidden,%interior hidden,
breakable,%before skip=10pt,after skip=10pt
}

\begin{document}
\section{Discrete Energy}
\begin{definition}
  The \( K \)-energy of a discrete point set \( \omega_N \) for some kernel \( K: A \times A \rightarrow \mathbb{R} \cup \{\infty\} \) for metric space \( (A, \rho) \) is
  \begin{equation}
  \label{eq:1}
  E_K(\omega_N) \coloneq \sum\limits_{i=1}^N \sum\limits_{\substack{j = 1 \\ j \neq i}}^N K(x_i, x_j) = \sum\limits_{i \neq j} K(x_i, x_{j}).
  \end{equation}
\end{definition}

\begin{definition}
  The minimal discrete \( N \)-point \( K \)-energy of A is
  \begin{equation}
  \label{eq:2}
  \mathcal{E}_K(A,N) \coloneq \inf \left\{ E_K(\omega_N) \mid \omega_N \subset A \right\}.
  \end{equation}
\end{definition}

Note that as \( N \rightarrow \infty \), \( \mathcal{E} \) grows at least as fast as \( N^2 \).

\begin{definition}
  The Riesz s-kernel is
  \begin{equation}
  \label{eq:3}
  K_s(x,y) \coloneq \begin{cases}
    \rho(x,y)^{-s}, & s\geq 0, \\
    -\rho(x,y)^{-s}, & s < 0
  \end{cases}.
  \end{equation}
\end{definition}

A discrete minimal Riesz energy function \( \mathcal{E}_s(A,N) = g(s) \) is continuous, for \( s > 0 \).

\begin{definition}
  The logarithmic kernel is how we define the \( s=0 \) kernel, and is
  \begin{equation}
  \label{eq:4}
  K_{\log}(x,y) \coloneq \log \frac{1}{\rho(x,y)},
  \end{equation}
  with associated \( E_{\log}(\omega_N) \) and \( \mathcal{E}_{\log}(A,N) \).
\end{definition}

Notice that minimizing the log kernel is maximizing the the product of distances.

\begin{definition}
  The Gaussian kernel is
  \begin{equation}
  \label{eq:5}
  G_t(x,y) \coloneq \exp(-t\rho(x,y)^2), \quad t > 0.
  \end{equation}
\end{definition}

Since \(\Gamma(s/2) = \int\limits_0^{\infty} e^{-t} t^{\frac{s}{2} -1} dt\) for \( s>0 \),
\begin{equation}
\label{eq:6}
K_s(x,y) = \frac{1}{\Gamma(s/2)} \int\limits_0^{\infty} G_t(x,y)t^{\frac{s}{2} -1} dt.
\end{equation}

\begin{definition}
  \( f: I \rightarrow \mathbb{R} \) is completely monotone on \( I \) if \(\forall n \geq 0\),
  \begin{equation}
  \label{eq:7}
  (-1)^n f^{(n)} \geq 0.
  \end{equation}
\end{definition}

\begin{thm}
  \( f \) is completely monotone \( (0, \infty) \) if and only if
  \begin{equation}
  \label{eq:8}
  f(x) = \int\limits_0^{\infty} e^{-xt} d\mu_f(t)
  \end{equation}
  for some positive Borel measure supported on \( [0,\infty] \).
\end{thm}

Hence, if \( f \) completely monotone, \( K(x,y) = f(\rho(x,y)^2) \) can be written as
\begin{equation}
\label{eq:9}
K(x,y) = \int\limits_0^{\infty} G_t(x,y) d\mu_f(t).
\end{equation}
\end{document}
